\section{System Design}

\Cref{fig:arch} shows the architecture. Each block maps to one or more requirements of \cref{tab:req}.

\begin{figure}[ht]
  \centering
  \begin{tikzpicture}[node distance=1.5cm and 1.4cm, every node/.style={font=\small},
      blk/.style={draw=atl-primary, fill=atl-light, rounded corners=3pt, minimum width=2.5cm, minimum height=0.95cm, thick, align=center},
      arr/.style={-{Stealth[length=2.2mm]}, thick, atl-secondary}]
    \node[blk] (s) {Sensor};
    \node[blk, right=of s] (m) {Microcontroller};
    \node[blk, right=of m] (u) {User interface};
    \node[blk, below=of m] (l) {Data logger};
    \draw[arr] (s) -- (m); \draw[arr] (m) -- (u); \draw[arr] (m) -- (l);
  \end{tikzpicture}
  \caption{System architecture.}\label{fig:arch}
\end{figure}

\subsection{Design Decisions}
\begin{enumerate}
  \item The sampling period is fixed by a hardware timer (R1).
  \item Data are stored in non-volatile memory (R2).
\end{enumerate}
